% The house layer of The Rhythm of Primes, after the GAI-MUA papers (00 campaign/styles/rs-header.tex): leading,
% colours, running heads, the title block, parts and headings, the theorem environments, the Lean anchors, the
% references list. Article form after the Annals of Mathematics: parts, sections numbered through the paper,
% statements numbered within sections, proofs ending with a square.
\linespread{1.06}
\setlength{\parskip}{0pt}
\setlength{\emergencystretch}{3em}      % an overfull line is a gate failure: stretch, never overrun
\clubpenalty=10000
\widowpenalty=10000
\displaywidowpenalty=10000
\setcounter{tocdepth}{1}     % the contents list parts and sections; the bookmarks keep the subsections
\AtBeginDocument{\urlstyle{same}}
\numberwithin{equation}{section}

\definecolor{rpteal}{HTML}{0B5F63}      % the darkest teal of the AURORA scale (7.0 : 1 on white)
\definecolor{rprule}{HTML}{A7C2C5}
\definecolor{rpgrey}{HTML}{4A575B}
\hypersetup{colorlinks=true,linkcolor=rpteal,citecolor=rpteal,urlcolor=rpteal,filecolor=rpteal,
  bookmarksdepth=3,pdfstartview=FitH}

% Running heads: the author on left-hand pages, the short title on right-hand pages, the folio outside.
\pagestyle{fancy}
\fancyhf{}
\renewcommand{\headrulewidth}{0pt}
\fancyhead[CE]{\footnotesize\scshape\RPauthor}
\fancyhead[CO]{\footnotesize\itshape\RPshorttitle}
\fancyhead[LE,RO]{\footnotesize\thepage}
\fancypagestyle{rpfirst}{\fancyhf{}\fancyfoot[C]{\footnotesize\thepage}}
\fancypagestyle{rppart}{\fancyhf{}}

\makeatletter
% amsart restores \topskip from \normaltopskip when it sets a page style (OpenShurik, 11-19): set both.
\setlength{\normaltopskip}{10pt plus 1pt}\setlength{\topskip}{\normaltopskip}

% ---- the title block ---------------------------------------------------------------------------------------------
% Above: version and date on the left, DOI (once assigned) and licence on the right, a hairline beneath. Then title,
% subtitle, author, one line with affiliation, e-mail and ORCID iD; then the abstract, MSC and keywords at 85 % of
% the measure.
\newcommand{\RPmaketitle}[1]{%
  \thispagestyle{rpfirst}%
  \begingroup\parindent\z@
  \vspace*{-6mm}%
  {\footnotesize\color{rpgrey}\RPversion\enspace\textperiodcentered\enspace\RPdate\hfill
   \ifx\RPdoi\@empty\else\href{https://doi.org/\RPdoi}{doi:\RPdoi}\enspace\textperiodcentered\enspace\fi
   \href{https://creativecommons.org/licenses/by/4.0/}{\RPlicence}\par}%
  \vspace{1.4mm}{\color{rprule}\hrule height .4pt}%
  \vspace{12mm}%
  \begin{center}
    {\fontsize{17.28}{21.5}\selectfont\bfseries\RPtitle\par}%
    \vspace{2.8mm}%
    {\fontsize{12}{15.2}\selectfont\itshape\RPsubtitle\par}%
    \vspace{6.5mm}%
    {\normalsize\scshape\RPauthor\par}%
    \vspace{1.4mm}%
    {\footnotesize\href{\RPaffiliationurl}{\RPaffiliation}\enspace\textperiodcentered\enspace
     \href{mailto:\RPemail}{\RPemail}\enspace\textperiodcentered\enspace\url{\RPorcid}\par}%
  \end{center}
  \vspace{5.5mm}%
  {\leftskip .075\textwidth \rightskip .075\textwidth
   \fontsize{9.5}{12.4}\selectfont{\scshape Abstract.}\enspace\ignorespaces #1\par
   \vspace{2.4mm}%
   \footnotesize\textit{2020 Mathematics Subject Classification.}\enspace\RPmsc.\par
   \vspace{.4mm}%
   \textit{Key words and phrases.}\enspace\RPkeywords.\par}%
  \vspace{7mm}%
  \endgroup}

% ---- parts and headings --------------------------------------------------------------------------------------------
% A part has a page of its own (layout gate L1): its number in small capitals, its name, and the summary that the
% part's introduction expands, between two short hairlines. Sections are numbered through the paper.
\newcounter{rppart}
\renewcommand{\therppart}{\Roman{rppart}}
\newcommand{\RPpart}[2]{% #1 label, #2 the name
  \clearpage\thispagestyle{rppart}%
  \refstepcounter{rppart}\label{#1}%
  \phantomsection\addcontentsline{toc}{part}{Part~\therppart.\texorpdfstring{\quad}{ } #2}%
  \vspace*{.28\textheight}%
  \begin{center}
    {\color{rprule}\rule{.28\textwidth}{.4pt}}\par\vspace{3.2mm}%
    {\small\scshape\color{rpteal}Part \therppart}\par\vspace{1.6mm}%
    {\fontsize{15}{18.5}\selectfont\bfseries #2\par}\vspace{3.4mm}%
    {\color{rprule}\rule{.28\textwidth}{.4pt}}\par
  \end{center}
  \clearpage}
\def\@secnumfont{\bfseries}
\def\section{\@startsection{section}{1}{\z@}{1.75\linespacing\@plus.4\linespacing\@minus.2\linespacing}%
  {.7\linespacing}{\normalfont\fontsize{12.5}{15}\selectfont\bfseries\raggedright}}
\def\subsection{\@startsection{subsection}{2}{\z@}{1.15\linespacing\@plus.3\linespacing\@minus.1\linespacing}%
  {.4\linespacing}{\normalfont\normalsize\bfseries\raggedright}}
\def\subsubsection{\@startsection{subsubsection}{3}{\z@}{.8\linespacing\@plus.3\linespacing}%
  {-.6em}{\normalfont\normalsize\itshape}}
% A heading keeps at least three lines of its text on its page (layout gate P2): a section asks for room for its
% heading and four lines, a subsection for its heading and three (with the space above them), or starts the next page.
\let\rp@section\section\let\rp@subsection\subsection
\renewcommand{\section}{\par\Needspace{8\baselineskip}\rp@section}
\renewcommand{\subsection}{\par\Needspace{7\baselineskip}\rp@subsection}
% The appendices open on a page of their own, like the parts (layout gate P0).
\newcommand{\RPappendices}{%
  \clearpage\thispagestyle{rppart}%
  \pdfbookmark[0]{Appendices}{appendices}%   % a bookmark only: the contents list the appendices themselves
  \addtocontents{toc}{\protect\addvspace{.6em}}%
  \vspace*{.28\textheight}%
  \begin{center}
    {\color{rprule}\rule{.28\textwidth}{.4pt}}\par\vspace{3.2mm}%
    {\fontsize{15}{18.5}\selectfont\bfseries Appendices\par}\vspace{3.4mm}%
    {\color{rprule}\rule{.28\textwidth}{.4pt}}\par
  \end{center}
  \clearpage}
% Contents: parts in small capitals, sections with their page numbers.
% The contents are set in the small size, as in the Annals, so that they end on a well-filled page.
\let\rp@tableofcontents\tableofcontents
\renewcommand{\tableofcontents}{{\small\linespread{1}\selectfont\rp@tableofcontents}}
\def\l@part#1#2{\addpenalty{-\@highpenalty}\vskip .6em\@plus.2em
  {\noindent\scshape #1\hfill\upshape #2\par}\nobreak}

% ---- statements ----------------------------------------------------------------------------------------------------
% Theorems, propositions, lemmas and corollaries in italics; definitions, remarks, notation and examples upright.
% One counter for all, within the section ("Theorem 7.3"). The main theorems of the introduction are lettered.
\newtheoremstyle{rpplain}{.9\baselineskip}{.9\baselineskip}{\itshape}{}{\bfseries}{.}{.5em}
  {\thmname{#1}\thmnumber{ #2}\thmnote{\ (\mdseries #3)}}
\newtheoremstyle{rpdefinition}{.9\baselineskip}{.9\baselineskip}{\normalfont}{}{\bfseries}{.}{.5em}
  {\thmname{#1}\thmnumber{ #2}\thmnote{\ (\mdseries #3)}}
\newtheoremstyle{rpremark}{.6\baselineskip}{.6\baselineskip}{\normalfont}{}{\itshape}{.}{.5em}
  {\thmname{#1}\thmnumber{ #2}\thmnote{\ (#3)}}
% Each environment that shares the counter has an alias counter (aliascnt), so that hyperref and
% cleveref see its own name ("Lemma 4.5"), not that of the shared counter.
\newcommand{\RPsharedthm}[2]{\newaliascnt{#1}{theorem}\newtheorem{#1}[#1]{#2}\aliascntresetthe{#1}}
\theoremstyle{rpplain}
\newtheorem{theorem}{Theorem}[section]
\RPsharedthm{proposition}{Proposition}
\RPsharedthm{lemma}{Lemma}
\RPsharedthm{corollary}{Corollary}
\newtheorem{mainthm}{Theorem}
\renewcommand{\themainthm}{\Alph{mainthm}}
\newtheorem*{theorem*}{Theorem}
\theoremstyle{rpdefinition}
\RPsharedthm{definition}{Definition}
\RPsharedthm{hypothesis}{Hypothesis}
\RPsharedthm{construction}{Construction}
\theoremstyle{rpremark}
\RPsharedthm{remark}{Remark}
\RPsharedthm{notation}{Notation}
\RPsharedthm{example}{Example}

% ---- Lean anchors --------------------------------------------------------------------------------------------------
% \lean{New.grh} (several names separated by commas) ties the statement in which it stands to the declarations of the
% Lean development that prove it. Nothing is printed in the text; the anchors are written, with the number and the
% page of the statement, to \jobname.lnk, from which Appendix C is set and against which the anchor and coverage
% gates check the rebuilt Lean environment.
% The anchor follows the statement, outside its group, so the number of the last statement is kept
% globally at the end of each statement environment (an equation inside it has ended by then).
\newwrite\rp@lnk
\AtBeginDocument{\immediate\openout\rp@lnk=\jobname.lnk\relax}
\def\rp@laststmt{}\def\rp@lastkind{}
% An anchor that stands in running text (a definition in a paragraph after the last statement) belongs to the
% subsection, not to that statement: a new paragraph after a statement clears the flag (LaTeX para hooks).
\newif\ifrp@afterstmt
\newcommand{\rp@track}[2]{\AddToHook{env/#1/end}{\xdef\rp@laststmt{\@currentlabel}\gdef\rp@lastkind{#2}%
  \global\rp@afterstmttrue}}
\AddToHook{para/begin}{\global\rp@afterstmtfalse}
\def\rp@here{\ifnum\c@subsection>\z@ \thesubsection\else\thesection\fi}
\rp@track{theorem}{Theorem}\rp@track{proposition}{Proposition}\rp@track{lemma}{Lemma}
\rp@track{corollary}{Corollary}\rp@track{mainthm}{Theorem}\rp@track{definition}{Definition}
\rp@track{hypothesis}{Hypothesis}\rp@track{construction}{Construction}\rp@track{remark}{Remark}
% \lean[eq:key]{names} anchors a numbered display (or a definition in running text) instead.
\newcommand{\lean}[2][]{\ifx\relax#1\relax
    \ifrp@afterstmt
      \protected@write\rp@lnk{}{\string\RPleanrow{\rp@lastkind}{\rp@laststmt}{\thepage}{\detokenize{#2}}}%
    \else
      \protected@write\rp@lnk{}{\string\RPleanrow{Text}{\rp@here}{\thepage}{\detokenize{#2}}}%
    \fi
  \else
    \protected@write\rp@lnk{}{\string\RPleanrowref{#1}{\thepage}{\detokenize{#2}}}%
  \fi\ignorespaces}
% Appendix C: one row per anchored statement; the names break after each comma. tools/build.py
% rewrites main.lnk into main.lnk.prev with the characters of Lean names made safe for pdfTeX.
\newcommand{\RPleanrow}[4]{#1~#2 & #3 & \RPleannames{#4}\\}
\newcommand{\RPleanrowref}[3]{Display~\eqref{#1} & #2 & \RPleannames{#3}\\}
\newcommand{\RPleanrowcont}[1]{ & & \RPleannames{#1}\\}   % continuation of a long row (tools/build.py)
\newcommand{\RPleannames}[1]{{\footnotesize\ttfamily\RP@names#1,\@nil}}
\def\RP@names#1,#2\@nil{#1\ifx\relax#2\relax\else,\allowbreak\ \RP@names#2\@nil\fi}
% The rows are read into a token list first: \input cannot stand inside an alignment.
\newtoks\rp@rows
\newread\rp@in
\newcommand{\RPleantable}{\IfFileExists{\jobname.lnk.prev}{%
  \global\rp@rows{}%
  \openin\rp@in=\jobname.lnk.prev\relax
  \begingroup\endlinechar=-1
  \loop\unless\ifeof\rp@in
    \read\rp@in to\rp@line
    \ifx\rp@line\@empty\else\global\rp@rows\expandafter{\the\expandafter\rp@rows\rp@line}\fi
  \repeat
  \endgroup\closein\rp@in
  \begin{longtable}{@{}l r >{\raggedright\arraybackslash}p{.66\textwidth}@{}}\toprule
  Statement & Page & Lean declarations\\\midrule\endhead
  \the\rp@rows\bottomrule\end{longtable}}{}}
% Appendix C as printed (D12): one row per section, read from \jobname.lnksum.prev, which tools/build.py writes from
% the anchors and the section labels of the previous pass; the full table above goes to out/supplement.
% A section: number, title, page, statements, declarations; then its module folders across the full measure.
\newcommand{\RPleansumrow}[5]{#1 & #2 & #3 & #4 & #5\\*}
\newcommand{\RPleansumfolders}[1]{\multicolumn{5}{@{\hspace{1.6em}}>{\raggedright\arraybackslash}p{\dimexpr\linewidth-1.6em\relax}@{}}{\footnotesize\ttfamily\hyphenchar\font=-1 #1}\\[1pt]}
\newcommand{\RPleansummary}{\IfFileExists{\jobname.lnksum.prev}{%
  \global\rp@rows{}%
  \openin\rp@in=\jobname.lnksum.prev\relax
  \begingroup\endlinechar=-1
  \loop\unless\ifeof\rp@in
    \read\rp@in to\rp@line
    \ifx\rp@line\@empty\else\global\rp@rows\expandafter{\the\expandafter\rp@rows\rp@line}\fi
  \repeat
  \endgroup\closein\rp@in
  \begin{longtable}{@{}l >{\raggedright\arraybackslash}p{.66\textwidth} r r r@{}}\toprule
  \S & Section and Lean module folders & Page & Stat. & Decl.\\\midrule\endhead
  \the\rp@rows\bottomrule\end{longtable}}{}}

% ---- references ----------------------------------------------------------------------------------------------------
% tools/refs.py writes generated/references.tex from ledger/refs.yaml, one \bibitem per verified entry.
\renewcommand{\bibliofont}{\small\raggedright}   % ragged, as in the GAI-MUA papers: DOIs and URLs do not stretch lines
\makeatother

\newcommand{\RPepigraph}[2]{%
  \begin{flushright}\begin{minipage}{0.78\textwidth}\raggedleft\itshape #1\par\smallskip
  \normalfont\small #2\end{minipage}\end{flushright}\bigskip}
\endinput
